\slajdid{09_3-s052}
\begin{frame}[label=09_3-s052]
\gusttekst\setlength{\abovedisplayskip}{3pt}\setlength{\belowdisplayskip}{3pt}\setlength{\abovedisplayshortskip}{2pt}\setlength{\belowdisplayshortskip}{2pt}Problem koji nastaje je šta ako je selekcioni signal $A=0$. U tom slučaju bi prolazni tranzistor bio zakočen i kapacitivnost na izlazu bi ostala na nivou na kojem je bila pre dovođenja selekcionog signala na 0 odnosno izlaz bi bio u stanju visoke impedanse. Nije nam to cilj. Zato moramo obezbediti i šta se dešava u ovom slučaju, odnosno vraćamo se na početnu ideju, selektorskih logičkih funkcija.
\begin{center}\begin{tikzpicture}[x=.85cm,y=.65cm,font=\sitneoznake,>=Latex]\ctikzset{bipoles/length=.6cm}\begin{scope}[shift={(0,1.6)}]\draw(-.35,0)--(.35,0);\draw(-.4,-.13)--(.4,-.13);\draw(-.3,0)--(-.3,.4)--(-1,.4)node[ocirc]{}node[left]{$B$};\draw(.3,0)--(.3,.4)--(1,.4);\draw(0,-.13)--(0,-.8)node[ocirc]{}node[below]{$A$};\node[above]at(0,.1){};\end{scope}\begin{scope}[shift={(0,-0.4)}]\draw(-.35,0)--(.35,0);\draw(-.4,-.13)--(.4,-.13);\draw(-.3,0)--(-.3,.4)--(-1,.4)node[ocirc]{}node[left]{$C$};\draw(.3,0)--(.3,.4)--(1,.4);\draw(0,-.13)--(0,-.8)node[ocirc]{}node[below]{$\overline A$};\node[above]at(0,.1){};\end{scope}\draw(1,2)--(1,0)--(1.8,0)node[ocirc]{}node[above]{Y}--(2.4,0);\draw(2,0)to[C,l=$C_L$](2,-1.5)node[ground]{};\end{tikzpicture}
\end{center}
\end{frame}
